\chapter{Higher-Order ODEs and Systems}

\section{Higher-order linear equations}

\begin{methodbox}{Constant coefficients, any order}
For $a_ny^{(n)} + \dots + a_1y' + a_0y = 0$, the characteristic equation is $a_nr^n + \dots + a_1r + a_0 = 0$. Each root gives solutions:
\begin{itemize}
  \item a real root $r$ of multiplicity $m$ gives $e^{rx},\ xe^{rx},\ \dots,\ x^{m-1}e^{rx}$;
  \item a complex pair $\alpha\pm i\beta$ of multiplicity $m$ gives $x^{k}e^{\alpha x}\cos\beta x$ and $x^{k}e^{\alpha x}\sin\beta x$ for $k = 0,\dots,m-1$.
\end{itemize}
The general solution is a sum of $n$ independent solutions with $n$ constants. An IVP needs $n$ initial conditions: $y(x_0), y'(x_0), \dots, y^{(n-1)}(x_0)$.
\end{methodbox}

\begin{examplebox}{Triple root}
Solve $y''' - 3y'' + 3y' - y = 0$. The characteristic equation is $(r-1)^3 = 0$, so $r=1$ three times:
\[ y = (c_1 + c_2x + c_3x^2)\,e^{x}. \]
\end{examplebox}

\begin{examplebox}{Fourth-order}
Solve $y^{(4)} - y = 0$. Then $r^4 = 1$, so $r = 1, -1, i, -i$:
\[ y = c_1e^{x} + c_2e^{-x} + c_3\cos x + c_4\sin x. \]
\end{examplebox}

\begin{examplebox}{Third-order IVP}
Solve $y''' - 6y'' + 11y' - 6y = 0$, \; $y(0)=0$, \; $y'(0)=1$, \; $y''(0)=0$.

$r^3 - 6r^2 + 11r - 6 = (r-1)(r-2)(r-3)$. So $y = c_1e^{x} + c_2e^{2x} + c_3e^{3x}$. The conditions give
\[ c_1 + c_2 + c_3 = 0,\quad c_1 + 2c_2 + 3c_3 = 1,\quad c_1 + 4c_2 + 9c_3 = 0. \]
Subtract the first from the second: $c_2 + 2c_3 = 1$. Subtract the second from the third: $2c_2 + 6c_3 = -1$. Then $2c_3 = -3$, so $c_3 = -\tfrac32$, $c_2 = 4$, $c_1 = -\tfrac52$:
\[ y = -\tfrac52e^{x} + 4e^{2x} - \tfrac32e^{3x}. \]
\end{examplebox}

For non-homogeneous higher-order equations, undetermined coefficients and variation of parameters work in the same way as in Chapter 2.

\section{Turning one ODE into a system}

Any $n$-th order ODE can be written as $n$ first-order equations. This is how computers solve ODEs (Chapter 5).

\begin{examplebox}{Conversion}
Write $y'' + 3y' + 2y = 0$ as a system. Let $u_1 = y$ and $u_2 = y'$. Then
\[ u_1' = u_2, \qquad u_2' = -2u_1 - 3u_2 , \]
or in matrix form $\mathbf{u}' = \begin{pmatrix}0 & 1\\ -2 & -3\end{pmatrix}\mathbf{u}$.
\end{examplebox}

\section{Linear systems: the eigenvalue method}

A \textbf{system} of linear ODEs with constant coefficients is written
\[ \mathbf{y}' = A\,\mathbf{y}, \qquad \mathbf{y}(0) = \mathbf{y}_0 , \]
where $A$ is a constant matrix. Such systems model coupled tanks, coupled circuits and coupled masses.

\begin{methodbox}{Eigenvalue method}
\begin{enumerate}
  \item Find the eigenvalues $\lambda$ from $\det(A - \lambda I) = 0$.
  \item For each $\lambda$, find an eigenvector $\mathbf{v}$ from $(A - \lambda I)\mathbf{v} = \mathbf{0}$.
  \item \textbf{Distinct real $\lambda$:} $\mathbf{y} = c_1\mathbf{v}_1e^{\lambda_1t} + c_2\mathbf{v}_2e^{\lambda_2t}$.
  \item \textbf{Complex pair $\lambda = \alpha \pm i\beta$:} take the real and imaginary parts of $\mathbf{v}e^{\lambda t}$ as two real solutions.
  \item \textbf{Repeated $\lambda$ with only one eigenvector $\mathbf{v}$:} find $\mathbf{w}$ from $(A-\lambda I)\mathbf{w} = \mathbf{v}$. Then $\mathbf{y} = c_1\mathbf{v}e^{\lambda t} + c_2(t\mathbf{v} + \mathbf{w})e^{\lambda t}$.
  \item Use the initial vector $\mathbf{y}(0)$ to find the constants.
\end{enumerate}
\end{methodbox}

\begin{derivbox}{Why the eigenvalue method works}
For $\mathbf{y}' = A\mathbf{y}$ try $\mathbf{y} = \mathbf{v}\,e^{\lambda t}$, where $\mathbf{v}$ is a constant vector. Then $\mathbf{y}' = \lambda\mathbf{v}e^{\lambda t}$ and $A\mathbf{y} = A\mathbf{v}\,e^{\lambda t}$. The equation holds if
\[ A\mathbf{v} = \lambda\mathbf{v} \iff (A - \lambda I)\mathbf{v} = \mathbf{0}. \]
A non-zero $\mathbf{v}$ exists only if $\det(A-\lambda I) = 0$. So $\lambda$ must be an eigenvalue and $\mathbf{v}$ an eigenvector. The system is linear, so sums of such solutions are solutions (superposition). Eigenvectors for different eigenvalues are independent, so for an $n\times n$ matrix with $n$ distinct eigenvalues we get $n$ independent solutions, which is a general solution.

\textbf{Complex eigenvalues.} If $A$ is real, complex eigenvalues come in conjugate pairs. The complex solution $\mathbf{v}e^{\lambda t}$ splits into a real part $\mathbf{u}$ and an imaginary part $\mathbf{w}$. Since $A$ is real, $\mathbf{u}'+i\mathbf{w}' = A\mathbf{u} + iA\mathbf{w}$, so the real and imaginary parts separately satisfy the system. Both are real solutions.

\textbf{Repeated eigenvalue with one eigenvector.} Try $\mathbf{y} = (t\mathbf{v} + \mathbf{w})e^{\lambda t}$. Then $\mathbf{y}' = (\mathbf{v} + \lambda t\mathbf{v} + \lambda\mathbf{w})e^{\lambda t}$ and $A\mathbf{y} = (tA\mathbf{v} + A\mathbf{w})e^{\lambda t} = (\lambda t\mathbf{v} + A\mathbf{w})e^{\lambda t}$. Setting these equal gives $\mathbf{v} + \lambda\mathbf{w} = A\mathbf{w}$, that is $(A - \lambda I)\mathbf{w} = \mathbf{v}$, as in the method box.
\end{derivbox}

\begin{examplebox}{Distinct real eigenvalues (IVP)}
Solve $x' = x + 2y$, \; $y' = 2x + y$, \; $x(0)=2$, \; $y(0)=0$.

$A = \begin{pmatrix}1&2\\2&1\end{pmatrix}$. The characteristic equation $(1-\lambda)^2 - 4 = 0$ gives $\lambda = 3, -1$.

For $\lambda=3$: $(A-3I)\mathbf{v} = \begin{pmatrix}-2&2\\2&-2\end{pmatrix}\mathbf{v}=\mathbf 0$, so $\mathbf{v}_1 = (1,1)^T$.
For $\lambda=-1$: $\begin{pmatrix}2&2\\2&2\end{pmatrix}\mathbf{v} = \mathbf 0$, so $\mathbf{v}_2 = (1,-1)^T$.

\[ \begin{pmatrix}x\\y\end{pmatrix} = c_1\begin{pmatrix}1\\1\end{pmatrix}e^{3t} + c_2\begin{pmatrix}1\\-1\end{pmatrix}e^{-t}. \]
At $t=0$: $c_1 + c_2 = 2$ and $c_1 - c_2 = 0$, so $c_1 = c_2 = 1$:
\[ x = e^{3t} + e^{-t}, \qquad y = e^{3t} - e^{-t}. \]
\textit{Check:} $x' = 3e^{3t} - e^{-t}$ and $x + 2y = 3e^{3t} - e^{-t}$. \checkmark
\end{examplebox}

\begin{examplebox}{Complex eigenvalues}
Solve $\mathbf{y}' = \begin{pmatrix}-1 & 2\\ -2 & -1\end{pmatrix}\mathbf{y}$.

$(-1-\lambda)^2 + 4 = 0$ gives $\lambda = -1 \pm 2i$. For $\lambda = -1+2i$: $\begin{pmatrix}-2i & 2\\ -2 & -2i\end{pmatrix}\mathbf v = \mathbf 0$. The first row gives $v_2 = i\,v_1$, so $\mathbf{v} = (1, i)^T$.

Now $\mathbf{v}e^{\lambda t} = e^{-t}\begin{pmatrix}1\\ i\end{pmatrix}(\cos 2t + i\sin 2t) = e^{-t}\begin{pmatrix}\cos 2t\\ -\sin 2t\end{pmatrix} + i\,e^{-t}\begin{pmatrix}\sin 2t\\ \cos 2t\end{pmatrix}$. The real and imaginary parts are two real solutions:
\[ x = e^{-t}(c_1\cos 2t + c_2\sin 2t), \qquad y = e^{-t}(-c_1\sin 2t + c_2\cos 2t). \]
The solution spirals into the origin because the real part $-1$ is negative.
\end{examplebox}

\begin{examplebox}{Repeated eigenvalue: two tanks in series}
Tank A (100 L) holds 10 kg of salt. Tank B (100 L) holds pure water. Fresh water flows into A at 5 L/min, A flows into B at 5 L/min, and B flows out at 5 L/min. Let $x$ and $y$ be the salt in A and B. Then
\[ x' = -0.05x, \qquad y' = 0.05x - 0.05y, \qquad x(0)=10,\ y(0)=0 . \]
The matrix is $\begin{pmatrix}-0.05&0\\0.05&-0.05\end{pmatrix}$ with a repeated eigenvalue $-0.05$. Solve in order:
$x = 10e^{-0.05t}$. Then $y' + 0.05y = 0.5e^{-0.05t}$, and with integrating factor $e^{0.05t}$ we get $(ye^{0.05t})' = 0.5$, so $y = (0.5t + C)e^{-0.05t}$. Since $y(0)=0$, $C=0$:
\[ y = 0.5\,t\,e^{-0.05t}. \]
$y$ is largest at $t = 20$ min, where $y = 10e^{-1}\approx 3.68$ kg. (The factor $t$ in the solution is the sign of a repeated eigenvalue.)
\end{examplebox}

\begin{practicebox}
\begin{enumerate}
  \item Find the general solution of $y''' - 4y' = 0$.
  \item Find the general solution of $y^{(4)} + 2y'' + y = 0$.
  \item Find the general solution of $x' = 4x + y$, \; $y' = 2x + 3y$.
  \item Solve the system in problem 3 with $x(0) = 3$, $y(0)=0$.
  \item Write $y''' + 2y'' - y' + 3y = \sin t$ as a first-order system.
\end{enumerate}
\end{practicebox}

\begin{answerbox}
(1) $r = 0, 2, -2$: $y = c_1 + c_2e^{2x} + c_3e^{-2x}$. \quad (2) $(r^2+1)^2 = 0$: $y = (c_1 + c_2x)\cos x + (c_3 + c_4x)\sin x$. \quad (3) $\lambda^2 - 7\lambda + 10 = 0$, $\lambda = 5, 2$; eigenvectors $(1,1)$ and $(1,-2)$: $x = c_1e^{5t} + c_2e^{2t}$, $y = c_1e^{5t} - 2c_2e^{2t}$. \quad (4) $c_1 + c_2 = 3$, $c_1 - 2c_2 = 0$ give $c_1 = 2$, $c_2 = 1$: $x = 2e^{5t} + e^{2t}$, $y = 2e^{5t} - 2e^{2t}$. \quad (5) With $u_1=y$, $u_2=y'$, $u_3=y''$: $u_1' = u_2$, $u_2' = u_3$, $u_3' = -3u_1 + u_2 - 2u_3 + \sin t$.
\end{answerbox}
